\section{Introduction}
\label{Introduction}
\setlength{\parindent}{20pt}
\justifying

\subsection{\large{\emph{Statistical learning}}}
\label{SL}
\noindent
When presented with absolutely novel input, our mind attempts to make sense of it as quickly as possible. This can occur without explicit control, and begins with the automatic and unconscious detection of the underlying statistics within the input. Slowly, over time, the accumulation of lower-order statistical information can be leveraged to form higher-order inferences and conscious knowledge of what we are perceiving. Within the statistical learning community, debate has been ongoing about the quality of statistical learning: whether it is modality-specific or domain-general, whether it is one process applied to different situations or is several different computations intervening at different stages of processing, whether it is aided by other cognitive abilities.

One of the most elegant accounts of the nature of statistical learning comes from \textcite{Thiessen2013}. They outline that information in novel input holds various statistical regularities, which may be used in different ways to understand the input. Conditional statistics map the frequency-related relationship between elements and determine to which extent one predicts the other. This is what was measured in \textcite{Saffran1996}'s pioneering experiment. The continuous speech stream ``bidakupadotigolabubidaku...'' is made of three-syllable ``words'': ``bidaku $\vert$ padoti $\vert$ golabu...''. Over the course of the experiment, the transitional probabilities (the probabilities of one element following another one; TPs) are computed: within-word syllables ``bi'', ``da'', and ``ku'' always follow each other (TP = 1), while between-word syllable TPs are much less reliable (``ku'' and ``pa'' would only have a TP of 0.33). In this way, conditional statistics tell participants that individual syllables form words, and that there are different words in this input. These kind of statistical analyses determine boundaries in stimuli, and allow for them to be stored as separate representations. This is the nature of the process the authors called extraction.

Other types of statistics exist. Distributional statistics, where elements can be grouped categorically according to a central tendency, is a type of learning called integration. This is what operates when doing phonological classification, for example: while there is variation in the detailed acoustic properties of a phoneme, listeners are able to classify all of these as variations around a same theme. Extraction and integration form a complementary system: input is divided into individual clusters through extraction, and grouped together thanks to integration.  

In the world of linguistic input, the goal is that of communication, of understanding the message behind the perceived randomness. The mind is highly motivated to learn informative cues that might aid the formation of lexical items. Many aspects of language echo across different existing languages, and the knowledge of other language structures can provide an approximate blueprint of what to search for in other languages. In linguistics, morphology holds one of the keys to decrypting a word. Should a reader not know the meaning of the rare English word ``irenic'', the ability to detect the suffix ``-ic'' within it allows us to conclude that this word is an adjective. Morphemes hold information on part of speech, structure, and semantics of a word. Affixes, through being much higher in frequency than stems, are more familiar to us, and allow us to understand new words. Affix detection and tracking can therefore allow for easier word learning. This is supported by a statistical learning study testing learning of affixation patterns in an artificial lexicon by \textcite{Lelonkiewicz2020}. Through purely passive exposure to these artificial affixed words, participants were found to be sensitive to the presence of a morpheme, and even more so if it was placed in its learned correct position within the word.


\subsection{\large{\emph{Bilingual statistical learning}}}
\label{BiSL}
\noindent
%% Bi env
Similar statistical learning paradigms have also been developed in bilingual experimental designs. However, this raises additional complications. In a natural bilingual setting, the content of speech is not the only thing that needs to be tracked, but also the language in which the message is communicated. For a bilingual to converse in one language, this may also entail the inhibition of others, and the usage of regularities only belonging to the active language. When speaking with other bilinguals, the active language might change unexpectedly. This code switching may be more fluent or frequent depending on speaker habits and contextual constraints. This requires a listener to continuously track which language is being spoken and detect information indicating a change of language. These cues then need to lead to quick shift in internal language contexts (\cite{Antovich2018}). 

The structural sensitivity hypothesis developed in \textcite{Kuo2012} posits that bilinguals may be more efficient at the identification of these language-specific cues and of new regularities in linguistic input. This is a crucial ability, as even monolingual input can be rife with variability. How, then, can one determine whether a change in inflection represents a change in language, or merely an accent? There could be two possible sources of variability: variability within the current structure, or variability that indicates a change of structure (\cite{Bulgarelli2016}; \cite{Bulgarelli2018}). When mapping the pattern of regularities they are familiar with, bilinguals therefore need to form multiple representations for their knowledge of the statistics of each of their languages. These, however, may overlap. It could be that bilinguals deploy different strategies to resolve this conflict. Over the course of learning and with much exposure, a bilingual's statistical representations may become so fine-tuned as to allow even very small variations to act as a cue of a different language (\cite{Antovich2020}).

%% Bi segmentation
With the importance of language membership cues in mind, some bilingual variations of \textcite{Saffran1996}'s experiment have been utilised to examine bilingual segmentation abilities. In \textcite{Antovich2018}'s bilingual auditory statistical learning design, words similar to that used in the \textcite{Saffran1996} study belonged to one of two languages. One language was spoken entirely by a male voice, and the other by a female voice, and this acted as the language membership cue. Fourteen-month-old infants were exposed to interleaved segments from language 1 and 2. Analysis revealed that only the bilingual infants were able to segment words in both languages. Bilingual infants may be better equipped to track more complex statistical regularities, and/or might do so more efficiently than monolinguals. In a 2019 literature review by \textcite{Weiss2020}, research from similar experiments indicates that, for interleaved statistical learning, regularities are typically only picked up when they are paired with a clear indexical cue. The risk of not having a clear clue in sequentially-presented languages is to obtain a primacy effect, in which learners only show segmentation abilities in the first language they were given (\cite{Bulgarelli2016}).

Another interesting study by \textcite{Bartolotti2011} examined how similarity of the language membership cues might complicate bilingual learning. In this experiment, artificial Morse languages were presented to participants, with multiple possible segmentation cues present. When these interfered with each other to a high degree, performance for all participants was around chance. However, when there was low interference between cues, bilinguals outperformed monolinguals. Beyond this, the bilinguals with more bilingual experience were better at the task than those with less experience. More bilingual experience might therefore aid in the process of integrating segmentation cues, or to inhibit one of the cues and learn just thanks to the other. This conclusion is similar to that from \textcite{Weiss2020}, which proposes that prior language experience modulates expectations and the processing of novel artificial languages. Early bilingual experience is thought to have a role in developing inhibitory control, information encoding, the generalisation of memory. Inhibitory control, in particular, could facilitate language context changes (\cite{Antovich2018}), and reduce the cognitive load of processing input in two languages concurrently.

%% Bi grammar learning
However, simple segmentation in a laboratory is far from the level of complexity that assails bilinguals in their usual environment (\cite{Frost2019}). Bilingual grammar learning tasks, while few, provide a more difficult task for bilinguals in the lab. One early such study by \textcite{Nation1986} found that multilinguals performing an implicit grammar learning task had fewer misses in their responses, and tended to respond more intuitively. Similar ideas were brought up in a paper by \textcite{Roncaglia2022} on machine classification of bilingual and multilingual performance in a grammar task: multilinguals seemed to perform with more metalinguistic awareness.

Past research therefore seems to suggest a bilingual, or even multilingual advantage for statistical learning tasks that are themselves designed to have a bilingual component. They seem to be able to track patterns that are more complex, and perform some tasks more efficiently.


\subsection{\large{\emph{Limitations to the current literature}}}
\label{LitLimitations}
\noindent
%% Def of a bilingual participant
However, some important methodological points can be made regarding the analysis of bilingual data. In a great many experiments, participants are themselves chunked as ``monolinguals'' or ``bilinguals''. These labels are very reductive, when the bilingual experience is much richer than the number of languages spoken fluently. \textcite{Bassetti2022} declare that studies with bilingual populations should report age of acquisition, proficiency, use, dominance, knowledge of other languages, linguistic and cultural contexts, type of exposure, language of schooling, and length of residence in countries speaking their languages, for a start. This is especially important since it is increasingly difficult to find ``true'' monolingual participants these days. Especially in countries where English is not the primary cultural language, a lot of that language is picked up online, in films, or on social media. The minimum would be to record their language experience according to the aforementioned metrics to check that they do fulfil the desired ``monolingual'' characteristics for an experiment. The reality is that these metrics fall on a continuum of multilingualism. Any researcher would be hard-pressed to conclusively determine at which level of second-language experience someone can be deemed a bilingual, and nor would a researcher want to separate two bilingual profiles that sit close together at such a boundary. We therefore need to be using continuous measures to chart multilingual participants.

%% Discrepancies in how different metrics are computed (bilingual dominance/balance)
Differences also exist in how some multilingual metrics are defined and computed. This is notably the case for the measure of multilingual balance. This is applied in two different ways: by thinking about language dominance (which is the ``main'' language) and language balance (the extent to which a bilingual's languages are of equal importance and ability). These are two sides of the same coin, but are not always discussed or computed as such. Language dominance has, in previous studies, been calculated as the difference between a participant's global language scores (\cite{Gertken2014}; \cite{Onnis2018}), or as the ratio of scores (\cite{Flege2002}; \cite{Sheng2014}). Language balance has sometimes been looked at as the difference only of \emph{proficiency} in various languages (\cite{Bartolotti2017}). Further complications arise when considering the study of multilinguals, which requires the comparison of more than two numbers. The only solution to this issue that this author has found is that of \textcite{Gullifer2019}, who use the notion of entropy to compare multiple language scores to determine balance.

%% Bilinguals or multilinguals
As is perhaps evident from the above sections, the majority of literature on bilingual statistical learning stops at the bilingual participant. However, a similar argument may be used in the determination of bilingualism. How much knowledge of a third language would cause a bilingual to be instead called a trilingual? Our multilingual world is full to the brim with as many different and unique multilingual profiles as there are multilinguals. A reductive approach not only ignores most of that complexity, but also introduces many confounds. As researchers, we may be missing fascinating effects that live purely in this world of messy multilingualism if we stop looking after two languages (\cite{Bassetti2022}). To give research every chance of picking up on these, we need to develop and standardise analytical methods that work across the multilingual spectrum.

%% Explicitness of the cue
In past bilingual statistical learning paradigms, the language membership cue only systematically supported learning (at least for bilinguals) when this was very overt. Some studies also explicitly told participants that there were two languages in the artificial input (cite{Bulgarelli2016}), a context in which success may be easier to find (\cite{Gebhart2009}). Very perceptually salient cues may also facilitate learning (\cite{Weiss2009}). However, language membership cues such as speaker voice may not have been interpreted by participants as entirely reliable: the same person may speak several languages when interacting with a fellow bilingual. Statistical learning can take place completely unconsciously in the absence of salient cues. It is therefore conceivable that a paradigm can be designed with some form of bilingual statistical learning happening thanks to an entirely implicit, entirely reliable language membership cue, voiding any top-down interferences and prior knowledge.


\subsection{\large{\emph{The current study}}}
\label{CurrentStudy}
\noindent
The current study tested the use of a completely implicit cue to language membership, that of morphological co-occurrence patterns that do not overlap across two languages. Repeated exposure to artificial words with these patterns should allow for participants to sort the presented words into two separate languages. Participants speaking any number of languages were accepted, and a variety of continuous measures of multilingualism were examined in conjunction with test results to evaluate both the success of the paradigm, and the usefulness of the individual metrics in capturing multilingual profiles.

















%\subsection{\large{\emph{Statistical Learning}}}
%\label{SL}

%\noindent
%Defined broadly, statistical learning (SL) is the process by which the brain unconsciously and automatically picks up regularities in one's environment. With reference to languages, this process is used to understand grammar rules, the structure of sentences, and also of individual words. A classic study paradigm examining the brain's ability to understand the latter is the segmentation task by \textcite{Saffran1996}. A stream of continuous speech composed of syllables is given to participants to listen to (for example, ``bidakupalabugobidakulabu''). With more exposure to the stream, each syllable is understood as an individual unit (``bi\textbar da\textbar ku\textbar pa\textbar la\textbar bu\textbar go\textbar bi\textbar da\textbar ku\textbar la\textbar bu''), and participants come to implicitly understand that some of these units are always presented together (``\textbf{bi\textbar da\textbar ku}\textbar pa\textbar la\textbar bu\textbar go\textbar \textbf{bi\textbar da\textbar ku}\textbar la\textbar bu''). This is due to the tracking of \textbf{transitional properties} (TPs), the frequency linking one part to another, with items with high-frequency TPs being clustered together.

%The extraction/integration framework (\cite{Thiessen2013}) provides a useful structure for these SL processes. To begin, several different kinds of statistical regularities are defined. \textbf{Conditional statistics} allow us to track co-occurrences in order to uncover relationships between stimuli. \textbf{Distributional statistics}, on the other hand, are used for the identification of an aspect common to a subset of stimuli, mentally clustering stimuli around a central prototype. Finally, and of particular interest in the current study, is \textbf{cue learning}. This is when perceptible cues are associated with certain stimulus properties in order to support learning. However, Thiessen and colleagues also outline two distinct mental processes in SL. \textbf{Extraction} is the clustering of stimuli into separate chunks and their storage as distinct representations, while \textbf{integration} is the grouping of separate stimuli as guided by the underlying statistics. According to this, Saffran-like segmentation tasks expose participants to conditional statistics, which require extraction to be parsed into chunks.

%Morphology, or the grammatical structure of words, holds an important place in SL research. Indeed, many languages are densely populated not only with stems (the core part of a word) but also affixes, which modify the stem to form a new word (the addition of the suffix ``ing'' transforms the verb stem ``sing'' into the gerund-participle ``sing-ing''). While studies like \textcite{Saffran1996} combine syllables into meaningless chunks, studies involving deliberate morphemic structures marry statistical learning and linguistics more closely, and allow for conclusions to be drawn about how we process more naturalistic linguistic stimuli. However, few SL studies extend in this direction. Of these, the paper by \textcite{Lelonkiewicz2020} has had the greatest impact on the design of the current study. Lelonkiewicz and colleagues expose participants to stimuli in an artificial lexicon, each word stimulus of an affix + random stem form. In the subsequent testing phase, they found that participants attributed new stimuli as coherent with the previous set of words if these contained one of the familiarised affixes. What's more, this effect was stronger if the affix was in the correct position (affix + stem if it was in the first position during familiarisation, stem + affix if it was in the second position). The consequences of these findings is that participants segmented the input and tracked these repeating affix-like chunks, and were sensitive their position within a word. In other words, they treated these chunks like real affixes, and applied grammatical rules when processing them: an affix can be combined with lots of different stems (``sing-ing'', but also ``laugh-ing'', ``talk-ing'', ``eat-ing'', etc.), as long as their position within the word is preserved (``ing'' can only be fitted \emph{after} a stem).


%\subsection{\large{\emph{Bilingual Statistical Learning}}}
%\noindent
%This paradigm can also be used to study the interplay of multiple language systems, by giving participants stimuli from two different languages. An additional element needs to be introduced in these studies: a \textbf{language membership cue} indicating which system each stimulus belongs to. \textcite{Antovich2018} proposed a paradigm whereby this language cue is the gender of the speaker's voice. For example, where a male voice is in green and a female voice is in red, participants would hear ``\textcolor{mygreen}{bidaku}\textcolor{myred}{palabugo}\textcolor{mygreen}{bidaku}\textcolor{myred}{labu}''. The clear pitch difference between a male and female voice enables participants to segment and track co-occurrences in both languages simultaneously. There is a growing body of research examining the effect of participants' language experience on their performance in these bilingual tasks. While this field needs more evidence to concretise any conclusions, it seems that bilinguals globally are able to perform better than monolinguals, especially in situations in which the task is more complex, or the cue is less explicit.

%\subsection{\large{\emph{Limitations to the Current Literature}}}
%\label{Limitations}
%\noindent \underline{Definition of a Bilingual Participant}

%\noindent
%People with different linguistic backgrounds have so many dimensions on which they differ, from their environment to their food, that for the sake of a clean study, the approach of limiting the data sample to one population has been adopted. However, aiming for more generalisation of the results, it is necessary to expand participation beyond a set number of languages spoken or place lived. 

%Many studies propose simple categorisations of their participants into mono- or bilinguals. However, multilingual research diverges in opinions on which element(s) of multilingualism contribute more to different cognitive processes. This suggests that multilingualism cannot be reduced down to a single dimension along which participants can vary, and that categorising participants in a discrete way eliminates a lot of the individual variability that might be driving specific effects. Indeed, a participant who has grown up as an early bilingual might not be equivalent to one who has been monolingual for most of their life. Similarly, someone using two languages every day might perform differently from a bilingual having regularly use of only one of their languages. The erasure of this richness describing multilingual experience in past studies limits results to very general findings about bilingual populations, without delving into what aspect of multilingualism might be driving performance.

%\noindent \underline{Bilinguals or Multilinguals}

%\noindent
%Furthermore, past research has principally focused on monolinguals and bilinguals. However, in a constantly more connected world, there is no clear reason why to stop at bilinguals. If multilingualism is to be explored on a continuous scale along continuous variables (such as Age of Acquisition of languages, percent of daily use of each language, proficiency in languages, to name a few), then extending research to more multilingual populations will expand the amount of individual variability along these axes, and allow for clearer relationships between performance and multilingualism to emerge.

%\noindent \underline{Expicitness of the Cue}

%\noindent
%Ecological validity

%\subsection{\large{\emph{The Current Research}}}
%\label{Current_Study}

%\noindent
%The present study aims to present participants with a statistical learning task in a bilingual setting with an implicit language membership cue. This cue could only be inferred through the learning of the combinatorial structure underlying the words participants were presented with. The experiment therefore consisted in learning combinations of artificial morphemes. Participants of varied linguistic backgrounds were invited to participate in the study and complete a thorough multilingual experience questionnaire, the answers to which were distilled into several theory-driven and data-driven measures of multilingualism.